\documentclass[11pt]{article}

\usepackage[T1]{fontenc}
\usepackage[utf8]{inputenc}
\usepackage[margin=1in]{geometry}
\usepackage{amsmath}
\usepackage{amssymb}
\usepackage{booktabs}
\usepackage{array}
\usepackage{hyperref}
\usepackage{xcolor}

\hypersetup{
  colorlinks=true,
  linkcolor=blue,
  citecolor=blue,
  urlcolor=blue
}

\newcommand{\pshort}{P20}
\newcommand{\pname}{P2.0 rotary state-output}
\newcommand{\pcontrol}{P2.0 rotary state-output control}
\newcommand{\pctrlshort}{P2.0 control}
\newcommand{\preposhort}{P20-control}
\newcommand{\parcae}{Parcae}
\newcommand{\bx}{B(x)}
\newcommand{\tokens}{\mathrm{tokens}}

\title{\pcontrol{} as an Efficient Depth Substitute in Tiny Causal Language Models}
\author{Fractal Research Notes}
\date{April 16, 2026}

\begin{document}
\maketitle

\begin{abstract}
We report a bounded H100 proof ladder for a small recurrent control primitive,
\pcontrol{} (repository shorthand: \preposhort{}), inserted into a looped
middle-depth causal language-model scaffold.
The goal is not to claim that the primitive generally beats Transformers.
The narrower question is whether the primitive provides useful depth-like
computation under a fixed tiny-model experimental contract. In a 9.87M
parameter model trained for 491.52M OpenLLaMA-tokenized FineWeb tokens,
\pcontrol{} beats the strongest non-\pshort{} looped control across three seeds,
survives a speed-compensated control check, and remains ahead after a bounded
learning-rate sweep of non-\pshort{} looped baselines. Against pure attention,
\pcontrol{} beats 9L, 10L, and 12L attention controls on the seed-42 depth
screen and falls into a three-seed tie band with 14L attention while using
about 10\% fewer parameters. A 16L attention model, however, beats \pcontrol{}
on seed 42 while using more parameters, more memory, and lower throughput.
These results support a modest but useful claim: in this tiny-LM regime,
\pcontrol{} behaves like an efficient depth substitute, but it has not yet
shown a general advantage over sufficiently deeper attention. We close with a
proposed 30M-50M parameter experiment needed to determine whether the effect
survives beyond tiny scale.
\end{abstract}

\section{Motivation}

Transformer decoders remain the default architecture for language modeling
because causal self-attention provides exact content-based token interaction
and scales well with modern GPU kernels \cite{vaswani2017attention,
dao2022flashattention}. Recurrent and state-space alternatives, including
selective state-space models such as Mamba \cite{gu2023mamba}, offer a
different tradeoff: they can provide recurrent state updates and efficient
sequence processing, but they must compete with a very strong and highly
optimized attention baseline.

The experiment reported here asks a narrower question than ``can a recurrent
primitive replace attention?'' We instead test whether a recurrent primitive
can act as a control mechanism for repeated middle-depth computation inside a
causal decoder. The primitive is the \pname{} recurrent primitive, abbreviated
as \pshort{} in the repository because it is the runtime-frozen `P2.0' member
of the broader P2 predictive-core family. It is not inserted as a full
attention replacement. It controls the injection path of a looped middle-depth
scaffold, where a small set of middle attention blocks is reused for multiple
recurrent passes.

The practical motivation is simple. If a small control primitive lets a model
reach the quality of a deeper pure-attention model with fewer parameters or
similar wall-clock cost, then it may provide a useful architectural axis for
small and mid-scale language models. If the effect disappears when pure
attention is tuned or deepened, then the result is still useful: it defines the
boundary at which ordinary depth catches the recurrent-control scaffold.

\section{Architecture}

\subsection{Base model}

All runs use a causal decoder language model trained for next-token prediction.
The base tiny model has:

\begin{itemize}
  \item vocabulary size 32,000 using an OpenLLaMA SentencePiece tokenizer,
  \item model width $d=128$,
  \item four attention heads,
  \item local context length 256,
  \item feed-forward multiplier 4,
  \item eight physical transformer blocks unless otherwise stated.
\end{itemize}

The pure attention baseline is an ordinary stack of local causal attention
blocks. The \parcae{} family preserves the same language-model objective and
tokenization, but changes the execution scaffold.

\subsection{Looped middle-depth scaffold}

The \parcae{} scaffold splits the physical transformer stack into a prelude, a
middle recurrent segment, and a coda. For the default 8-layer model, the split
is:

\[
  3\ \text{prelude blocks}
  \rightarrow
  2\ \text{middle blocks looped twice}
  \rightarrow
  3\ \text{coda blocks}.
\]

The scaffold maintains a recurrent middle state. At each loop, it combines a
bounded injection from the prelude stream with a decayed previous state, applies
the middle attention blocks, and updates the recurrent state through a bounded
nonlinear delta. This is a stable looped-compute scaffold, not a faithful
reproduction of any external looped language-model paper.

\subsection{P2.0 rotary state-output primitive}

The name \pshort{} is internal shorthand, not a public architectural
description. The underlying primitive is a \pname{} recurrent scan. It descends
from the project's earlier contractive recurrent primitive line, but adds the
two contracts that the simpler P1-style update lacked: a transformed carry path
before the update and an emitted output that is distinct from the latent state.

For an input token representation $x_t$ and prior latent state $s_{t-1}$, the
primitive computes packed input projections
\[
  (a_t,\theta_t,c_t,r_t) = W_{\mathrm{in}} x_t
\]
and applies a gated transformed-state update
\[
  u_t = \operatorname{Rot}(W_s s_{t-1}, \theta_t),
\]
\[
  s_t = g(a_t) \odot u_t + (1 - g(a_t)) \odot \tanh(c_t),
\]
\[
  o_t = g(r_t) \odot s_t.
\]
Here $g(\cdot)$ is the bounded gate used by the primitive implementation,
$W_s$ is the recurrent state transform, $\operatorname{Rot}$ applies the
input-conditioned pairwise rotary transform, $s_t$ is latent memory, and $o_t$
is the emitted output returned to the surrounding model. The fast frozen
runtime used in the earlier Path 1 work exposes $W_{\mathrm{in}}$ as one packed
projection and uses a block-diagonal recurrent state transform for the Triton
scan path.

In this report, \pcontrol{} means that this \pname{} scan is used as a control
producer for the looped scaffold. It should not be read as a claim that the
primitive replaces attention inside the reported model.

\subsection{Control variants}

We evaluate four relevant variants:

\begin{description}
  \item[Attention-only.] A standard local causal attention stack.
  \item[\parcae{} looped attention.] The looped middle scaffold with a shared,
    learned scalar/channel injection.
  \item[\parcae{} \bx{} looped attention.] The same scaffold with a learned
    value projection and gate for the injected prelude stream.
  \item[\parcae{} \pcontrol{} looped attention.] The same scaffold with a
    \pname{} recurrent scan over the prelude stream. The primitive output is
    normalized and used to produce the injected value and gate for the middle
    loop.
\end{description}

The \pcontrol{} path is intentionally not a standalone sequence mixer in this
experiment. It is a control primitive attached to a looped attention scaffold.
This distinction matters for the claim boundary.

\section{Experimental Contract}

The H100 proof ladder used:

\begin{itemize}
  \item Modal 1x NVIDIA H100 80GB HBM3,
  \item PyTorch 2.10.0 with CUDA 12.8,
  \item bf16 training,
  \item Triton primitive backend,
  \item OpenLLaMA-tokenized FineWeb CC-MAIN-2024-10 cache,
  \item 750,000,544 train tokens and 10,029,403 eval tokens,
  \item sequence length 256,
  \item batch size 64,
  \item 64 eval batches,
  \item deterministic seed and data-seed pairs.
\end{itemize}

Each 30,000-step lane sees
\[
  30{,}000 \times 64 \times 256 = 491{,}520{,}000
\]
training tokens, about 65.5\% of the train split. The runner enforces a
no-repeat cap, so these runs do not wrap the training cache.

\section{Results}

\subsection{Initial four-lane promotion}

Table~\ref{tab:promotion} shows the seed-42 promotion run. Plain looped
middle-depth accounts for most of the gain over the weak 8L attention baseline.
\bx{} improves the looped scaffold slightly. \pcontrol{} improves it further.

\begin{table}[h]
\centering
\caption{Seed-42 four-lane H100 promotion, 30k steps.}
\label{tab:promotion}
\begin{tabular}{lrrrr}
\toprule
Lane & Params & Final loss & tok/s & Peak MB \\
\midrule
8L attention-only & 9.78M & 5.9347 & 1,067,937 & 6,774.79 \\
\parcae{} looped attention & 9.78M & 4.4132 & 878,842 & 6,983.55 \\
\parcae{} \bx{} looped attention & 9.81M & 4.3842 & 866,260 & 6,991.98 \\
\parcae{} \pctrlshort{} looped attention & 9.87M & 4.3012 & 793,615 & 7,046.88 \\
\bottomrule
\end{tabular}
\end{table}

\subsection{Equal-step comparison against \bx{}}

The cleanest same-scaffold comparison is \bx{} versus \pcontrol{}. Table
\ref{tab:bx} shows the three-seed equal-step result. \pcontrol{} wins all
three seeds.

\begin{table}[h]
\centering
\caption{Equal-step finalist repeat, 30k steps.}
\label{tab:bx}
\begin{tabular}{rrrrrr}
\toprule
Seed & \bx{} loss & \pctrlshort{} loss & Delta & \bx{} tok/s & \pctrlshort{} tok/s \\
\midrule
42 & 4.3842 & 4.3012 & -0.0830 & 866,260 & 793,615 \\
43 & 4.4658 & 4.2972 & -0.1687 & 908,138 & 832,650 \\
44 & 4.3518 & 4.2965 & -0.0553 & 863,253 & 792,434 \\
\midrule
Mean & 4.4006 & 4.2983 & -0.1023 & 879,217 & 806,233 \\
\bottomrule
\end{tabular}
\end{table}

The mean loss delta is -0.1023 in favor of \pcontrol{}, corresponding to a
perplexity ratio of $\exp(-0.1023) \approx 0.9027$. In this scaffold, the
primitive is not merely adding noise or capacity; it provides a repeatable
control advantage over the strongest non-\pshort{} looped injection path.

\subsection{Speed-compensated \bx{} pressure check}

Because \pcontrol{} is slower than \bx{}, we also gave \bx{} additional steps
to spend its speed advantage. This is not an exact stopwatch-identical test,
because observed H100 throughput varies by run, but it is a useful practical
pressure check.

\begin{table}[h]
\centering
\caption{Speed-compensated \bx{} pressure check.}
\label{tab:bxspeed}
\begin{tabular}{rrrrrr}
\toprule
Seed & \bx{} steps & \bx{} loss & \pcontrol{} loss & Delta & \bx{} train min \\
\midrule
42 & 32,750 & 4.4223 & 4.3012 & -0.1211 & 9.39 \\
43 & 32,725 & 4.5352 & 4.2972 & -0.2380 & 9.45 \\
44 & 32,685 & 4.4055 & 4.2965 & -0.1090 & 10.33 \\
\midrule
Mean & -- & 4.4543 & 4.2983 & -0.1560 & 9.72 \\
\bottomrule
\end{tabular}
\end{table}

The faster \bx{} control did not catch \pcontrol{} when given extra steps. In
fact, the speed-compensated \bx{} runs were worse than their 30k-step
counterparts at this learning rate, suggesting that simply extending the
control lane at fixed LR was not enough.

\subsection{Tuned non-\pshort{} controls}

We next gave non-\pshort{} looped controls a bounded learning-rate sweep on seed 42:
3e-4, 6e-4, 1e-3, and 1.5e-3. The best non-\pshort{} setting was plain \parcae{}
looped attention at LR 6e-4. It still lost to \pcontrol{} across three seeds.

\begin{table}[h]
\centering
\caption{Seed-42 learning-rate sweep for non-\pshort{} looped controls.}
\label{tab:lrsweep}
\begin{tabular}{rrr}
\toprule
LR & \parcae{} looped loss & \bx{} looped loss \\
\midrule
3e-4 & 4.5403 & 4.5320 \\
6e-4 & 4.3712 & 4.3890 \\
1e-3 & 4.4132 & 4.3842 \\
1.5e-3 & 7.7104 & 5.7993 \\
\bottomrule
\end{tabular}
\end{table}

\begin{table}[h]
\centering
\caption{Best tuned non-\pshort{} control versus \pcontrol{}.}
\label{tab:tuned}
\begin{tabular}{rrrrr}
\toprule
Seed & Tuned looped loss & \pcontrol{} loss & Delta & Tuned tok/s \\
\midrule
42 & 4.3712 & 4.3012 & -0.0700 & 935,472 \\
43 & 4.3771 & 4.2972 & -0.0800 & 913,677 \\
44 & 4.3781 & 4.2965 & -0.0816 & 856,470 \\
\midrule
Mean & 4.3755 & 4.2983 & -0.0772 & 901,873 \\
\bottomrule
\end{tabular}
\end{table}

The mean perplexity ratio from the mean loss delta is
$\exp(-0.0772) \approx 0.9257$. \pcontrol{} therefore survives the first
tuned-control rung, although it remains slower than the tuned looped control.

\subsection{Pure attention depth boundary}

The strongest caveat comes from deeper pure attention. We screened attention
depth at the same width and context length, using LR 6e-4 for the deeper
attention controls. Table~\ref{tab:depth} shows seed 42.

\begin{table}[h]
\centering
\caption{Pure attention depth screen, seed 42.}
\label{tab:depth}
\begin{tabular}{lrrrr}
\toprule
Lane & Params & Final loss & tok/s & Train min \\
\midrule
8L attention, LR 1e-3 & 9.78M & 5.9347 & 1,067,937 & 7.67 \\
9L attention, LR 6e-4 & 9.98M & 4.4459 & 1,091,606 & 7.50 \\
10L attention, LR 6e-4 & 10.17M & 4.3649 & 981,974 & 8.34 \\
12L attention, LR 6e-4 & 10.57M & 4.3304 & 897,467 & 9.13 \\
14L attention, LR 6e-4 & 10.97M & 4.3044 & 773,922 & 10.59 \\
16L attention, LR 6e-4 & 11.36M & 4.2758 & 721,719 & 11.35 \\
\pcontrol{}, LR 1e-3 & 9.87M & 4.3012 & 793,615 & 10.32 \\
\bottomrule
\end{tabular}
\end{table}

This screen shows the boundary clearly. \pcontrol{} beats the 9L, 10L, and
12L attention controls. It is essentially tied with 14L attention and loses to
16L attention on seed 42. The 16L model is larger, slower, and uses more CUDA
memory, but it has better loss.

We repeated the 14L comparison across seeds because it is the boundary closest
to \pcontrol{}.

\begin{table}[h]
\centering
\caption{14L attention boundary repeat.}
\label{tab:14l}
\begin{tabular}{rrrrrr}
\toprule
Seed & 14L loss & \pctrlshort{} loss & Delta & 14L tok/s & \pctrlshort{} tok/s \\
\midrule
42 & 4.3044 & 4.3012 & -0.0032 & 773,922 & 793,615 \\
43 & 4.2960 & 4.2972 & +0.0011 & 772,807 & 832,650 \\
44 & 4.3061 & 4.2965 & -0.0097 & 847,395 & 792,434 \\
\midrule
Mean & 4.3022 & 4.2983 & -0.0039 & 798,041 & 806,233 \\
\bottomrule
\end{tabular}
\end{table}

The mean \pcontrol{} edge over 14L attention is only -0.0039 loss, a
perplexity ratio of about 0.9961. This should be interpreted as a tie band, not
a decisive quality win. The practical result is nevertheless meaningful:
\pcontrol{} reaches the 14L attention boundary with about 10\% fewer
parameters and similar throughput.

\subsection{Speed-spend pressure against 16L attention}

We also let \pcontrol{} spend some of its speed advantage against the 16L
attention model. A 33k-step run reached 4.2835 loss, and a 34k-step run reached
4.2812 loss. Both narrowed the gap but did not catch the 16L attention result
of 4.2758 on seed 42.

\section{Discussion}

The core result is best described as an efficient-depth signal. The
\pcontrol{} scaffold does not simply beat all attention. It also does not
collapse once attention is strengthened. Instead, it occupies a useful middle
ground:

\begin{itemize}
  \item It reliably beats same-scaffold non-\pshort{} controls.
  \item It beats parameter-near and compute-near attention controls.
  \item It sits in a tie band with a much deeper 14L attention model.
  \item It loses to 16L attention at this context length and token budget.
\end{itemize}

This is the first result in this line of work that looks like a credible
architectural role rather than a primitive leaderboard curiosity. The primitive
is not replacing attention. It is controlling repeated middle-depth computation
and acting as a compact substitute for some ordinary depth.

There are two immediate interpretations. The optimistic interpretation is that
\pcontrol{} supplies a recurrent control state that helps the looped middle
blocks refine hidden representations more efficiently than adding several
ordinary attention layers. The skeptical interpretation is that looped middle
depth itself provides most of the gain, and \pcontrol{} is a small regularizer
or gating improvement whose advantage may vanish at larger scale or longer
training. The current data favors the first interpretation within the tiny
contract, but it cannot settle the scaling question.

\section{Limitations}

The result should not be overstated.

\begin{itemize}
  \item The models are tiny: roughly 10M parameters.
  \item Context length is only 256 tokens.
  \item The attention baselines were tuned only lightly.
  \item The strongest pure attention model tested, 16L, beat \pcontrol{}.
  \item We do not yet have a 30M-50M scale result.
  \item We do not yet have long-context evidence.
  \item We have not shown generation quality, downstream task quality, or
    quantization behavior for this specific scaffold.
\end{itemize}

The supported claim is therefore:

\begin{quote}
In the current 10M-parameter, 256-context Path 1 benchmark, \pcontrol{} acts as
an efficient depth substitute inside a looped attention scaffold. It beats
same-scaffold controls and reaches the 14L pure-attention boundary with fewer
parameters, but it does not beat sufficiently deeper attention.
\end{quote}

\section{Next Experiment}

To determine whether this is more than a tiny-model artifact, the next
experiment should move to 30M-50M parameters. A practical first target is:

\begin{itemize}
  \item width $d=320$,
  \item 16 physical layers for the \pcontrol{} scaffold,
  \item 8 heads,
  \item roughly 40.8M parameters,
  \item 1.0B to 1.25B training tokens,
  \item fixed OpenLLaMA tokenizer and FineWeb-derived train/eval split.
\end{itemize}

The control set should include:

\begin{itemize}
  \item the scaled \pcontrol{} scaffold,
  \item a parameter-near attention baseline,
  \item a stronger depth-near attention baseline,
  \item optionally a looped non-\pshort{} control if budget allows.
\end{itemize}

One H100 is sufficient for this next rung. Multi-GPU training is not required
because the models fit comfortably on one H100 and the current harness is
single-GPU. If parallel compute is available, it should be used for independent
lanes or seeds rather than distributed training.

A bounded compute request of 50-100 H100-hours would cover:

\begin{itemize}
  \item token-cache preparation and smoke tests,
  \item one-seed 30M-50M screening,
  \item finalist promotion to a larger token budget,
  \item seed repeats only if the first screen remains alive.
\end{itemize}

\section{Conclusion}

\pcontrol{} is not yet a transformer replacement result. It is a credible
efficient-depth result. In a controlled tiny-LM H100 ladder, \pcontrol{}
improves a looped middle-depth attention scaffold beyond matched non-\pshort{}
controls, reaches the 14L attention boundary with fewer parameters, and remains
competitive until pure attention is made substantially deeper. The next
scientific question is whether this role survives at 30M-50M parameters and
longer context. If it does, the primitive deserves a more serious scaling-law
study. If it does not, the current work still clarifies where recurrent control
helps and where ordinary attention depth catches up.

\bibliographystyle{plain}
\bibliography{references}

\end{document}
